\chapter{Choosing a Sampler}
\label{ch:choosing-sampler}

\newthought{The sampler decides which parameter points} \Jarvis\ evaluates. This part of the book
has one focused section per method; this chapter helps you pick the right family and explains the
few things every sampler shares. Each later section gives the \emph{algorithm idea} and the
\emph{exact \textsc{yaml} format} for that method.

\section{A quick decision guide}
\label{sec:choose-guide}

Match your goal to a starting point (Table~\ref{tab:choose}). Most first runs should use
\sampler{Random} or \sampler{Bridson} to validate the model, then switch to a specialised method
once the workflow is stable.

\begin{table}[ht]
  \footnotesize
  \begin{tabular}{@{}p{0.46\linewidth}l@{}}
    \toprule
    \textbf{Your goal} & \textbf{Start with} \\
    \midrule
    Quick workflow / plumbing test            & \sampler{Random} \\
    Broad space-filling coverage              & \sampler{Bridson} (or \sampler{Grid}) \\
    Replay a fixed list of points             & \sampler{CSV} \\
    Posterior / Bayesian sampling             & \textsc{mcmc} family \\
    Evidence / model comparison               & \sampler{Dynesty}, \sampler{MultiNest} \\
    Global best-fit / optimisation            & \sampler{Diver} \\
    \textsc{ml}-guided exploration            & \sampler{DNN} \\
    Profile a nuisance parameter (inner step) & \sampler{Profile1D} \\
    \bottomrule
  \end{tabular}
  \caption{Goal-to-sampler quick guide.}
  \label{tab:choose}
\end{table}

\section{The shared shape}
\label{sec:choose-shared}

Every sampler is configured in the \yamlkey{Sampling} block (Chapter~\ref{ch:sampling-block}):
\yamlkey{Method} names it, \yamlkey{Variables} and \yamlkey{LogLikelihood} are the same
everywhere, and the method's own dials live in \yamlkey{Bounds} (or, for a few samplers, in a
method-named key). The skeleton is always:

\begin{yamlwin}
Sampling:
  Method: "SamplerName"
  Variables: [ ... ]
  Bounds: { ... }        # the dials for this sampler
  LogLikelihood: [ ... ]
\end{yamlwin}

\begin{jwarn}
Control keys are \emph{not} interchangeable between samplers. Even similar-sounding keys differ
in name and meaning---\yamlkey{Point number} belongs to \sampler{Random},
\yamlkey{Radius} to \sampler{Bridson}, \yamlkey{num\_chains} to the \textsc{mcmc} family. Always
use the section for your chosen method as the source of truth.
\end{jwarn}

\section{Main and nuisance samplers}
\label{sec:choose-main-nuisance}

\Jarvis\ distinguishes two roles:

\begin{description}[style=nextline, leftmargin=1.2em]
  \item[Main sampler] chooses the physics parameter points that drive the scan (everything in
        this chapter's table).
  \item[Nuisance sampler] an optional \emph{inner} step run around each main point to
        profile/optimise nuisance parameters. It augments---does not replace---the main sampler.
        \sampler{Profile1D} is the one provided (Chapter~\ref{ch:nuisance}).
\end{description}

\section{The sampler families at a glance}
\label{sec:choose-families}

Table~\ref{tab:sampler-map} maps every method to its family and primary control keys. The rest of
Part~\ref{part:samplers} is one section per method.

\begin{table*}[ht]
  \footnotesize
  \begin{tabular}{@{}p{1.0in}p{2.35in}p{2.45in}@{}}
    \toprule
    \textbf{Family (chapter)} & \textbf{Methods} & \textbf{Primary keys} \\
    \midrule
    Basic (Ch.~\ref{ch:basic-samplers})        & \sampler{Random}, \sampler{Grid}, \sampler{Bridson}, \sampler{CSV} & \yamlkey{Point number} / per-var \yamlkey{num} / \yamlkey{Radius} / \yamlkey{CSV} \\
    ML \& evolutionary (Ch.~\ref{ch:ml-evolutionary}) & \sampler{DNN}, \sampler{Diver}            & \yamlkey{Bounds} / \yamlkey{run} \\
    Nested (Ch.~\ref{ch:nested})               & \sampler{Dynesty}, \sampler{MultiNest}            & \yamlkey{nlive}, \yamlkey{rseed}, \yamlkey{run\_nested} \\
    MCMC I (Ch.~\ref{ch:mcmc-1})               & \sampler{MCMC}, \sampler{AMMCMC}, \sampler{RobustAM}, \sampler{DRAM} & \yamlkey{num\_chains}, \yamlkey{num\_iters}, \yamlkey{proposal\_scale} \\
    MCMC II (Ch.~\ref{ch:mcmc-2})              & \sampler{DEMCMC}, \sampler{DREAM}, \sampler{DREAMLite}, \sampler{EnsembleMCMC}, \sampler{PTEnsemble}, \sampler{PTMCMC} & (above) $+$ \yamlkey{exchange\_interval}, \yamlkey{proposal\_scales} \\
    MCMC III (Ch.~\ref{ch:mcmc-3})             & \sampler{SliceMCMC}, \sampler{ESS}                & \yamlkey{num\_chains}, \yamlkey{num\_iters}, \yamlkey{proposal\_scale} \\
    Gradient (Ch.~\ref{ch:gradient})           & \sampler{MALA}, \sampler{HMC}, \sampler{NUTS}     & (\emph{placeholder} engines) \\
    Nuisance (Ch.~\ref{ch:nuisance})           & \sampler{Profile1D}                               & \yamlkey{Sampling.Nuisance} \\
    Experimental (Ch.~\ref{ch:experimental})   & \sampler{RLTPMCMC}                                 & tempering $+$ \textsc{rl} control \\
    \bottomrule
  \end{tabular}
  \caption{Every sampler, by family and primary control keys.}
  \label{tab:sampler-map}
\end{table*}

\begin{jnote}
Registered aliases: \alias{ToyMCMC}{MCMC} and \alias{DREAM-lite}{DREAMLite}. Either name selects
the same sampler.
\end{jnote}

\section{A note on diagnostics}
\label{sec:choose-diagnostics}

What to check after a run depends on the family: coverage/spread for space-filling samplers;
trace plots, acceptance rate, autocorrelation, and effective sample size for \textsc{mcmc};
evidence stability and live-point behaviour for nested samplers; best-fit improvement and
repeatability for \sampler{Diver}. The per-method sections point out the most useful checks.
